%============================= abstract.tex================================
\chapter*{Abstract}%
%\addcontentsline{toc}{chapter}{\numberline{}Abstract}%
\addcontentsline{toc}{chapter}{Abstract}%

This research presents a novel approach to detecting epileptic seizures leveraging the strengths of Machine Learning (ML) and Deep Learning (DL) algorithms in EEG signals. Epileptic seizures are neurological events with distinctive features found in Electroencephalography (EEG) that lend considerable credibility to researchers. Machine Learning (ML) and Deep learning (DL) algorithms have emerged as powerful feature extraction and classification tools in EEG signal analysis. Many studies have converted the EEG signals into either images and /or calculated time-frequency domain features and performed classification. This study focuses on classifying time-series data representation of EEG signals with machine learning-based classifiers by tuning parameters and deep learning-based One-Dimensional Convolutional Neural Network (1D CNN) methods. The primary objective is not only to determine the optimal classifier but also to emphasize critical metrics such as sensitivity, precision, and accuracy, which are critical in medical investigations, particularly for the early detection of diseases and patient care optimization. The UCI Epileptic Seizure Recognition dataset used in this study consists of time-series data points extracted from the EEG signals. The dataset has been preprocessed and fed to the classifiers, namely Extreme Gradient Boosting (XGBoost), TabNet, Random Forest (RF), One Dimensional Convolutional Neural Network, and achieved encouraging accuracies of 98\%, 96\%, 98\%, and 99\%, respectively. The proposed 1D-CNN model performed better than other state-of-the-art models concerning accuracy, sensitivity, precision, and recall. \\


\thispagestyle{plain}
%=======================================================================

